% 核心观点：表示应保留区分动作任务后果所需的关系与因果执行上下文，
% 而不只是使专家动作MSE的上界更小。
% 叙述逻辑：信息保留/决策利用分解 -> 关系保留的条件性界 -> 语义与关系监督的对应。
% 正文集中比较条件与适用域；标量实例和伪标签误差分解留附录，检验协议见实验。
% 边界：语义是引导关系表示的手段，不另设动作预测风险或泛化收益命题；
% 关系界约束可达决策能力，不证明特定架构优于RGB基线。
% 任务后果代理仅为分析见证函数；几何辅助损失不等于价值监督或策略改进证明。
\subsection{Semantic Inputs and Task-Supervised Representations}
\label{sec:representationanalysis}
% 任务相关表示需要保留决定动作后果的关系与上下文；以下先区分信息损失与决策实现误差，再分析关系估计误差如何约束表示损失。
Task-relevant representations must preserve the relations and context needed to distinguish action outcomes.
We first separate information loss from decision-realization error, then bound representation loss through relation-estimation error.

% \subsubsection{表示信息与决策能力}
\subsubsection{Represented Information and Decision Capability}
% 固定参考策略 $\pi_0$ 及一个时刻--历史评价分布 $\mu$。
Fix a reference policy $\pi_0$ and a time--history evaluation distribution $\mu$.
% 动作规则 $a$ 对应的平均失败概率定义为
Define the mean failure probability associated with an action rule $a$ as
$\mathcal L_\mu(a)=\E_\mu[Q_t^0(x,a)]$.
% 它衡量替换当前动作后沿 $\pi_0$ 继续执行的平均失败概率。
This quantity measures the mean failure probability after replacing the current action and following $\pi_0$ thereafter.

% 含时间索引的完整因果历史生成信息 $\mathcal F_h=\sigma(\mathcal H_t)$。
The complete causal history, including the time index, generates the information $\mathcal F_h=\sigma(\mathcal H_t)$.
% 编码表示 $Z=\phi(\mathcal H_t)$ 生成的信息为
% Z在此表示目标指令的完整决策输入：除学习特征外，还包含目标时刻/槽位、所需计划缓存及因果调度记录。
The representation $Z=\phi(\mathcal H_t)$ denotes the complete input determining the target command: learned features together with the target time and slot, required cached plan information, and causal scheduling records.
% 这些量均由含运行记录的因果历史确定；确定性读出包括槽位选择及所用后处理。
These quantities are determined by the causal history including runtime records; the deterministic readout includes slot selection and any postprocessing used.
% Z止于待评价决策读出之前；反馈可条件于已有基础计划，但不将其自身的最终修正动作作为输入。
The representation is taken before the evaluated decision readout: feedback may condition on an existing base plan, but not on its own final corrected command.
It generates
$\mathcal U=\sigma(Z)\subseteq\mathcal F_h$.
% 基础策略与反馈分支的不同编码分别作为 $Z$ 评价。
The encodings of the base policy and the feedback branch are evaluated separately as $Z$.
% 对任意信息集合 $\mathcal F$，以 $\mathfrak A(\mathcal F)$ 表示所有
For any information set $\mathcal F$, let $\mathfrak A(\mathcal F)$ denote all
% $\mathcal F$-可测且取值于 $\mathcal A$ 的动作规则，
$\mathcal F$-measurable action rules taking values in $\mathcal A$.
% 仅使用这些信息所能达到的平均失败概率下确界为
The infimum of the mean failure probability attainable using only this information is
\begin{equation}
 \mathcal L_\mu^*(\mathcal F)
 =\inf_{a\in\mathfrak A(\mathcal F)}\mathcal L_\mu(a).
 \label{eq:taskinformationrisk}
\end{equation}
% 其中sigma表示记录可辨别的信息，F可测规则仅依赖该信息。
Here, $\sigma(\cdot)$ denotes the information generated by the records, and an $\mathcal F$-measurable rule uses only that information.
% 各类动作规则在相同动作集合下比较，当前动作之后均沿同一参考策略执行。
The action-rule classes are compared over the same action set, with the same reference policy followed after the current action.

% 基于表示的实际动作规则 $\pi_Z$ 同时受到编码信息与决策机制的限制。
An implemented action rule $\pi_Z\in\mathfrak A(\mathcal U)$ is constrained both by the encoded information and by its decision mechanism.
% 将编码丢失信息所造成的最优平均失败概率差记为表示损失
Define the gap in optimal mean failure probability due to information lost in encoding as the representation loss
$\Delta_{\mathrm{rep}}(\mathcal U)
=\mathcal L_\mu^*(\mathcal U)-\mathcal L_\mu^*(\mathcal F_h)$,
% 将实际决策相对同一表示下最优规则的差距记为决策实现误差
and the gap between the implemented rule and the optimum under the same representation as the decision-realization error
$\epsilon_{\mathrm{dec}}(\pi_Z)
% =\mathcal L_\mu(\pi_Z)-\mathcal L_\mu^*(\mathcal U)$，有
=\mathcal L_\mu(\pi_Z)-\mathcal L_\mu^*(\mathcal U)$. Then
\begin{equation}
 \mathcal L_\mu(\pi_Z)-\mathcal L_\mu^*(\mathcal F_h)
 =\Delta_{\mathrm{rep}}(\mathcal U)+\epsilon_{\mathrm{dec}}(\pi_Z).
 \label{eq:taskrepresentationdecomposition}
\end{equation}
% 该分解区分\emph{信息是否被保留}与\emph{保留的信息是否被有效利用}。
This decomposition distinguishes \emph{whether information is preserved} from \emph{whether the preserved information is effectively used}.
% 两项均以沿参考策略继续执行的平均失败概率计量且非负；$\epsilon_{\mathrm{dec}}$ 包含模型容量、残差限幅、
Both terms are nonnegative and measured in mean failure probability under reference-policy continuation. The term $\epsilon_{\mathrm{dec}}$ includes limitations due to model capacity, residual bounds,
% 有限数据与优化，以及动作监督和任务目标之间的差异。
finite data and optimization, and the mismatch between action supervision and the task objective.

% \subsubsection{关系误差何时影响任务决策}
\subsubsection{When Relation Errors Affect Task Decisions}
% 关系估计的准确性需要通过动作后果与任务决策联系起来。
The accuracy of relation estimates must be connected to task decisions through action outcomes.
% 给定可获得的历史，对未知物理状态取条件平均后的失败概率记为
Given the available history, denote the failure probability conditionally averaged over unknown physical states by
$\overline Q_\mu(\mathcal H_t,a)=\E_\mu[Q_t^0(x,a)\mid\mathcal F_h]$,
% 其中 $a$ 为固定的待比较动作，后续均沿参考策略执行。
where $a$ is the fixed action being compared, followed in each case by the reference policy.

% 关系变量 $G\in\R^{d_G}$ 描述相关对象的几何或语义关系，
The relation variable $G\in\R^{d_G}$ describes geometric or semantic relations among relevant objects,
% 因果上下文 $\chi$ 保留决策所需的本体信息与执行证据。
and the causal context $\chi$ retains proprioceptive information and execution evidence needed for decisions.
% chi的组成须在任务实例中预先限定，不能通过把完整历史放入chi使假设恒成立；未解释因素由epsilon_rel承担。
The contents of $\chi$ must be specified for the task instance, rather than set to the entire history to make the assumption automatic; unexplained effects remain in $\epsilon_{\mathrm{rel}}$.
% 表示 $Z$ 可读出关系估计 $\widehat G\in\R^{d_G}$ 与上下文 $\chi$，
The representation $Z$ admits readouts of a relation estimate $\widehat G\in\R^{d_G}$ and context $\chi$,
% 关系误差为 $\delta_G=\E_\mu\norm{G-\widehat G}$。
with relation error $\delta_G=\E_\mu\norm{G-\widehat G}$.
% 以下用分析函数 $f(G,\chi,a)$ 近似上述条件平均失败概率，
We use an analytical function $f(G,\chi,a)$ to approximate the conditional mean failure probability above.
% 其中d_G为关系维数。
Here, $d_G$ is the relation dimension.

% \begin{assumption}[关系条件下的任务后果近似]
\begin{assumption}[Relation-based approximation of action outcomes]
\label{ass:taskrelations}
% 在固定分布 $\mu$ 下，存在可测函数 $f$ 及常数 $\epsilon_{\mathrm{rel}},K_G\ge0$，使
Under a fixed distribution $\mu$, there exist a measurable function $f$ and constants $\epsilon_{\mathrm{rel}},K_G\ge0$ such that
\begin{align}
 \E_\mu\sup_{a\in\mathcal A}
 |\overline Q_\mu(\mathcal H_t,a)-f(G,\chi,a)|
 &\le\epsilon_{\mathrm{rel}},\nonumber\\
 |f(G,\chi,a)-f(\widehat G,\chi,a)|
 &\le K_G\norm{G-\widehat G},\quad a\in\mathcal A .
 \label{eq:taskrelationsassumption}
\end{align}
% 第二个条件在真实与估计关系的联合支持上几乎处处成立，
The second condition holds almost everywhere on the joint support of the true and estimated relations,
% 且 $\delta_G<\infty$；条件期望、上确界和动作选择满足附录~\ref{app:representationregularity}的正则条件。
and $\delta_G<\infty$. The conditional expectations, suprema, and action selections satisfy the regularity conditions in Appendix~\ref{app:representationregularity}.
\end{assumption}
% 其中，$\epsilon_{\mathrm{rel}}$ 为关系和上下文不足以解释任务后果的近似误差，
Here, $\epsilon_{\mathrm{rel}}$ is the approximation error due to task effects not captured by the relations and context,
% $K_G$ 为该后果对关系误差的敏感性常数，$\delta_G$ 为关系估计误差。
$K_G$ is the sensitivity constant with respect to relation errors, and $\delta_G$ is the relation-estimation error.
% 例如，相同几何下的接触或滑脱差异须由上下文保留，否则计入近似误差。
For example, contact and slip at the same geometry must be distinguished by the context or accounted for in the approximation error.

% \begin{proposition}[任务后果意义下的表示损失]
\begin{proposition}[Representation loss in task outcomes]
\label{prop:taskrepresentation}
% 在假设~\ref{ass:taskrelations}下，
Under Assumption~\ref{ass:taskrelations},
\begin{align}
 \Delta_{\mathrm{rep}}(\mathcal U)
 &\le 2(\epsilon_{\mathrm{rel}}+K_G\delta_G),\nonumber\\
 \mathcal L_\mu(\pi_Z)-\mathcal L_\mu^*(\mathcal F_h)
 &\le 2(\epsilon_{\mathrm{rel}}+K_G\delta_G)
       +\epsilon_{\mathrm{dec}}(\pi_Z).
 \label{eq:taskrelationbound}
\end{align}
\end{proposition}
% 证明见附录~\ref{app:representationproofs}。
The proof is given in Appendix~\ref{app:representationproofs}.
% RGB和增强编码器比较必须固定mu、pi0、动作集合与尺度、G定义与尺度、上下文及执行输入。
Comparisons between RGB and enhanced encoders use the same $\mu$, reference continuation $\pi_0$, action set and scale, relation definition and normalization, and context and execution inputs.
% 分别训练的表示并不必然信息嵌套，须分别检验假设与误差；仅当共享近似误差和敏感性上界时才能把更小delta解释为更紧上界。
Separately trained representations need not have nested information sets: the assumption and errors are evaluated for each, and smaller $\delta_G$ gives a tighter common bound only with shared bounds on $\epsilon_{\mathrm{rel}}$ and $K_G$.
% 该比较不从关系误差大小直接推出实际策略排序。
Relation-readout accuracy alone does not rank their implemented policies.
% 附录的理想化末步纠偏例子展示非零、有限关系敏感性下的信息价值，并非实际接触动力学模型。
An idealized final-step correction in Appendix~\ref{app:relationexample} illustrates the value of relations under nonzero, finite sensitivity; it is not a model of the actual contact dynamics.
% 关系因遮挡而未定义时，关系界仅适用于有效条件分布，无效情形仍计入总体任务结果。
When occlusion leaves relations undefined, the relation bound applies only to the valid conditional distribution, while invalid cases remain part of the overall task results.

% 关系读出与伪标签误差均需独立核查，不能以训练拟合代替理论中的评价误差。
Relation-readout and pseudo-label errors require independent assessment (Appendix~\ref{app:pseudolabel}); training fit cannot substitute for the evaluation errors used in the analysis.

% 上述结果表明，在相应条件下，准确保留任务相关关系与上下文可以约束表示损失；实际决策表现还取决于策略对这些信息的利用。
Under the stated conditions, accurately preserving task-relevant relations and context in $Z$ bounds representation loss; actual decision performance also depends on how the policy uses this information.
% 下一节进一步讨论执行中新增证据的利用。
The next section examines the use of additional evidence acquired during execution.
